\section{Resultados e Discussões}

Para as grandezas medidas repetidamente, adotou-se como resultado a média e, como incerteza, o desvio-padrão amostral das medidas,
\[
s=\sqrt{\frac{1}{N-1}\sum_{i=1}^{N}\left(q_i-\bar{q}\right)^2}.
\]

\subsection{Oscilação no ar e constante elástica da mola}

A massa do corpo foi medida três vezes, obtendo-se $10{,}90$, $10{,}91$ e $10{,}70\un{g}$, de modo que
\[
m=\frac{10{,}90+10{,}91+10{,}70}{3}=(10{,}84\pm0{,}12)\un{g}.
\]

A Tabela~\ref{tab:ar} apresenta os tempos de 10 oscilações medidos no ar, os períodos $T=t_{10}/10$, as frequências angulares obtidas pela Equação~\eqref{eq:wT} e as constantes elásticas obtidas pela Equação~\eqref{eq:k}, calculadas com a massa média.

\begin{table}[H]
\centering
\caption{Período, frequência angular e constante elástica para a oscilação no ar.}
\label{tab:ar}
\begin{tabular}{ccccc}
\toprule
Medida & $t_{10}$ (s) & $T_0$ (s) & $\omega_0$ (rad/s) & $k$ (N/m) \\
\midrule
1 & 13,91 & 1,391 & 4,517 & 0,2211 \\
2 & 13,99 & 1,399 & 4,491 & 0,2186 \\
3 & 14,04 & 1,404 & 4,475 & 0,2170 \\
4 & 13,97 & 1,397 & 4,498 & 0,2192 \\
5 & 14,02 & 1,402 & 4,482 & 0,2177 \\
\midrule
Média & --- & 1,399 & 4,493 & 0,2187 \\
Desvio-padrão & --- & 0,005 & 0,016 & 0,0016 \\
\bottomrule
\end{tabular}
\fonte{Elaborada pelos autores.}
\end{table}

Para a medida 1, por exemplo,
\[
\omega_0=\frac{2\pi}{1{,}391}=4{,}517\un{rad/s},
\qquad
k=\frac{4\pi^2\left(10{,}84\times10^{-3}\right)}{(1{,}391)^2}=0{,}2211\un{N/m}.
\]
Assim, o período de oscilação no ar, a frequência angular natural e a constante elástica da mola são
\[
T_0=(1{,}399\pm0{,}005)\un{s},
\qquad
\omega_0=(4{,}493\pm0{,}016)\un{rad/s},
\qquad
k=(0{,}219\pm0{,}002)\un{N/m}.
\]

\subsection{Oscilação na água: análise do período}

A Tabela~\ref{tab:agua} mostra os tempos de 10 oscilações, os períodos e as frequências angulares obtidos com o corpo oscilando na água.

\begin{table}[H]
\centering
\caption{Período e frequência angular para a oscilação na água.}
\label{tab:agua}
\begin{tabular}{cccc}
\toprule
Medida & $t_{10}$ (s) & $T_1$ (s) & $\omega_1$ (rad/s) \\
\midrule
1 & 14,60 & 1,460 & 4,304 \\
2 & 14,24 & 1,424 & 4,412 \\
3 & 14,13 & 1,413 & 4,447 \\
4 & 14,36 & 1,436 & 4,376 \\
5 & 14,27 & 1,427 & 4,403 \\
\midrule
Média & --- & 1,432 & 4,388 \\
Desvio-padrão & --- & 0,018 & 0,054 \\
\bottomrule
\end{tabular}
\fonte{Elaborada pelos autores.}
\end{table}

Logo,
\[
T_1=(1{,}432\pm0{,}018)\un{s},
\qquad
\omega_1=(4{,}388\pm0{,}054)\un{rad/s}.
\]

Comparando com o valor no ar, a frequência diminuiu $0{,}105\un{rad/s}$, ou seja, $2{,}3\%$. Os intervalos $\omega_0=[4{,}477;\,4{,}509]\un{rad/s}$ e $\omega_1=[4{,}334;\,4{,}442]\un{rad/s}$ não se sobrepõem, de modo que se pode afirmar que as frequências são diferentes. O sentido da diferença ($\omega_1<\omega_0$) é coerente com a Equação~\eqref{eq:w1}, que prevê frequência menor na presença de amortecimento. Pela Equação~\eqref{eq:gw}, o fator de amortecimento correspondente seria
\[
\gamma=\sqrt{(4{,}493)^2-(4{,}388)^2}=\sqrt{20{,}183-19{,}257}=0{,}96\un{rad/s},
\]
valor comparado, na próxima seção, com o obtido pelo decaimento das amplitudes.

\subsection{Oscilação na água: análise da variação de amplitude}

Com $T_1=1{,}432\un{s}$, os instantes de medida $t_n=n\,T_1/2$ são $0$; $0{,}716$; $1{,}432$; $2{,}148$; $2{,}864$; $3{,}580$ e $4{,}296\un{s}$. A Tabela~\ref{tab:decaimento} apresenta as amplitudes $|x_n|=|y_n-y_{eq}|$ medidas nas duas realizações do experimento, que resultaram em valores iguais, as amplitudes normalizadas por $x_0=15{,}0\un{cm}$ e seus logaritmos.

\begin{table}[H]
\centering
\caption{Amplitude de oscilação na água em função do tempo.}
\label{tab:decaimento}
\begin{tabular}{cccccc}
\toprule
$n$ & $t_n$ (s) & $|x_n|$ -- 1.ª medida (cm) & $|x_n|$ -- 2.ª medida (cm) & $|x_n/x_0|$ & $\ln|x_n/x_0|$ \\
\midrule
0 & 0,000 & 15,0 & 15,0 & 1,000 & $0{,}000$ \\
1 & 0,716 & 10,0 & 10,0 & 0,667 & $-0{,}405$ \\
2 & 1,432 & 7,0 & 7,0 & 0,467 & $-0{,}762$ \\
3 & 2,148 & 5,5 & 5,5 & 0,367 & $-1{,}003$ \\
4 & 2,864 & 4,5 & 4,5 & 0,300 & $-1{,}204$ \\
5 & 3,580 & 3,5 & 3,5 & 0,233 & $-1{,}455$ \\
6 & 4,296 & 3,0 & 3,0 & 0,200 & $-1{,}609$ \\
\bottomrule
\end{tabular}
\fonte{Elaborada pelos autores.}
\end{table}

As amplitudes decrescem progressivamente, e $\ln|x_n/x_0|$ diminui de forma aproximadamente linear com $t_n$, o que indica um decaimento exponencial, consistente com as Equações~\eqref{eq:xamort} e \eqref{eq:ln}. A regressão linear de $\ln|x_n/x_0|$ em função de $t_n$ resultou em coeficiente angular $-0{,}368\un{rad/s}$, com desvio-padrão $0{,}025\un{rad/s}$, coeficiente linear $-0{,}13$ e coeficiente de correlação $r=-0{,}988$. Portanto,
\[
\gamma=(0{,}37\pm0{,}03)\un{rad/s}.
\]
O coeficiente linear, que pela Equação~\eqref{eq:ln} deveria ser nulo, e a leve curvatura dos dados, com decaimento mais rápido no início do movimento, indicam que o modelo exponencial é apenas aproximado.

O decaimento das amplitudes mostra também que a energia mecânica do oscilador não se conserva. Nos pontos extremos a velocidade é nula e a energia é apenas potencial elástica, $E_n=\frac{1}{2}k\,x_n^2=\frac{1}{2}k\,x_0^2e^{-2\gamma t_n}$, de modo que a energia decai exponencialmente, sendo dissipada pelo atrito viscoso. Ao final de $3T_1$, por exemplo, $(x_6/x_0)^2=(0{,}200)^2=0{,}04$, isto é, restavam apenas $4\%$ da energia inicial.

Para verificar a consistência com a Equação~\eqref{eq:w1}, calculou-se o valor previsto de $\omega_1$ com o $\gamma$ obtido do decaimento:
\[
\omega_1=\sqrt{(4{,}493)^2-(0{,}37)^2}=4{,}477\un{rad/s},
\]
valor fora do intervalo medido, $\omega_1=(4{,}388\pm0{,}054)\un{rad/s}$. De forma equivalente, o valor $\gamma=0{,}96\un{rad/s}$, obtido na seção anterior a partir de $\omega_0$ e $\omega_1$, é cerca de 2,6 vezes maior que o obtido pelo decaimento. Portanto, os resultados não são estritamente consistentes com a Equação~\eqref{eq:w1}. Algumas causas plausíveis são:
\begin{enumerate}
    \item a frequência $\omega_0$ foi medida no ar, mas, na água, o corpo arrasta fluido consigo, o que aumenta sua massa efetiva e reduz a frequência natural no meio; assim, a Equação~\eqref{eq:gw}, aplicada com o $\omega_0$ do ar, atribui ao amortecimento uma redução de frequência que se deve, em parte, à massa efetiva, superestimando $\gamma$;
    \item o valor de $\gamma$ obtido pela Equação~\eqref{eq:gw} resulta da diferença entre dois números muito próximos ($20{,}18$ e $19{,}26\un{rad^2/s^2}$), o que o torna muito sensível a pequenos erros de cronometragem;
    \item a curvatura observada nos dados da Tabela~\ref{tab:decaimento} sugere que a força de atrito não é puramente proporcional à velocidade, de modo que o ajuste por uma única exponencial é aproximado;
    \item as amplitudes foram lidas visualmente na régua com a massa em movimento e nos instantes $t_n$ estimados pelo operador.
\end{enumerate}

\subsection{Oscilação forçada no ar}

Para cada rotação do motor, a frequência de excitação foi calculada por $\Omega=2\pi\cdot7/t_7$. Observa-se que os valores $7/t_7$ anotados no laboratório correspondem a rotações por segundo (Hz), sendo necessária a multiplicação por $2\pi$ para obter $\Omega$ em rad/s. A Tabela~\ref{tab:forcada_ar} reúne os dados obtidos.

\begin{table}[H]
\centering
\caption{Amplitude máxima da oscilação forçada no ar em função da frequência de excitação.}
\label{tab:forcada_ar}
\begin{tabular}{ccccc}
\toprule
Ponto & $t_7$ (s) & $7/t_7$ (Hz) & $\Omega$ (rad/s) & $x_0$ (cm) \\
\midrule
1 & 19,16 & 0,365 & 2,296 & 10,7 \\
2 & 12,56 & 0,557 & 3,502 & 21,4 \\
3 & 10,46 & 0,669 & 4,205 & 49,0 \\
R & 9,73 & 0,719 & 4,520 & $>49{,}0$ (não medida) \\
4 & 6,79 & 1,031 & 6,478 & 9,8 \\
5 & 4,03 & 1,737 & 10,914 & 8,3 \\
6 & 3,63 & 1,928 & 12,116 & 3,3 \\
\bottomrule
\end{tabular}
\fonte{Elaborada pelos autores.}
\end{table}

A amplitude cresce à medida que $\Omega$ se aproxima da ressonância e decai rapidamente acima dela. Na condição de ressonância, em que a amplitude ficou grande demais para ser medida, mediram-se $t_7=9{,}74\un{s}$ e $t_7=9{,}72\un{s}$, que correspondem a
\[
\Omega=\frac{2\pi\cdot7}{9{,}74}=4{,}516\un{rad/s}
\qquad\text{e}\qquad
\Omega=\frac{2\pi\cdot7}{9{,}72}=4{,}525\un{rad/s},
\]
de modo que a frequência de ressonância no ar é
\[
\Omega_R^{ar}=(4{,}520\pm0{,}007)\un{rad/s}.
\]

Comparando com a frequência do oscilador livre, $\omega_0=(4{,}493\pm0{,}016)\un{rad/s}$, a diferença é de $0{,}028\un{rad/s}$, ou seja, $0{,}6\%$. Os valores são, portanto, muito próximos, como esperado para um meio de baixo amortecimento, em que a Equação~\eqref{eq:Wr} fornece $\Omega_r\approx\omega_0$. Estritamente, a Equação~\eqref{eq:Wr} prevê $\Omega_r\le\omega_0$, enquanto o resultado experimental ficou ligeiramente acima, praticamente no limite dos intervalos de incerteza. Essa pequena diferença pode ser atribuída ao tempo de reação do operador ao iniciar e parar o cronômetro e à dificuldade de identificar visualmente a rotação exata em que a amplitude é máxima. Mesmo que o amortecimento no ar fosse tão grande quanto o medido na água ($\gamma=0{,}37\un{rad/s}$), a Equação~\eqref{eq:Wr} preveria $\Omega_r=4{,}462\un{rad/s}$, uma redução de apenas $0{,}7\%$ em relação a $\omega_0$.

\subsection{Oscilação forçada na água}

De modo análogo, os dados obtidos com o corpo na água estão na Tabela~\ref{tab:forcada_agua}.

\begin{table}[H]
\centering
\caption{Amplitude máxima da oscilação forçada na água em função da frequência de excitação.}
\label{tab:forcada_agua}
\begin{tabular}{ccccc}
\toprule
Ponto & $t_7$ (s) & $7/t_7$ (Hz) & $\Omega$ (rad/s) & $x_0$ (cm) \\
\midrule
1 & 22,30 & 0,314 & 1,972 & 7,5 \\
2 & 13,34 & 0,525 & 3,297 & 12,7 \\
3 & 11,63 & 0,602 & 3,782 & 19,0 \\
R & 10,445 & 0,670 & 4,211 & 22,9 \\
4 & 8,67 & 0,807 & 5,073 & 14,7 \\
5 & 7,77 & 0,901 & 5,661 & 8,2 \\
6 & 6,02 & 1,163 & 7,306 & 4,2 \\
\bottomrule
\end{tabular}
\fonte{Elaborada pelos autores.}
\end{table}

Na ressonância, mediram-se $t_7=10{,}52\un{s}$ e $t_7=10{,}37\un{s}$, que correspondem a $\Omega=4{,}181\un{rad/s}$ e $\Omega=4{,}241\un{rad/s}$. Assim,
\[
\Omega_R^{ag}=(4{,}21\pm0{,}04)\un{rad/s},
\]
com amplitude máxima $x_0=22{,}9\un{cm}$.

Comparando as Tabelas~\ref{tab:forcada_ar} e \ref{tab:forcada_agua}, nota-se que, no ar, a amplitude cresce muito perto da ressonância: já em $\Omega=4{,}205\un{rad/s}$ obteve-se $x_0=49{,}0\un{cm}$, e na ressonância a amplitude ultrapassou a faixa mensurável. Na água, o máximo foi de apenas $22{,}9\un{cm}$, e a amplitude varia de forma mais suave em torno de $\Omega_R$. Isso é coerente com a Equação~\eqref{eq:x0W}: quanto maior o fator de amortecimento, menor a amplitude máxima e mais larga a curva de ressonância. Quanto à posição do máximo, $\Omega_R$ diminuiu de $4{,}520\un{rad/s}$ no ar para $4{,}21\un{rad/s}$ na água, uma redução de $6{,}8\%$, também no sentido previsto pela Equação~\eqref{eq:Wr}, em que o aumento de $\gamma$ desloca a ressonância para frequências menores.

A consistência quantitativa com a Equação~\eqref{eq:Wr} depende do valor de $\gamma$ adotado, como mostra a Tabela~\ref{tab:comparacao}.

\begin{table}[H]
\centering
\caption{Valores de $\Omega_r$ na água previstos pela Equação~\eqref{eq:Wr}, com $\omega_0=4{,}493\un{rad/s}$.}
\label{tab:comparacao}
\begin{tabular}{lcc}
\toprule
Origem de $\gamma$ & $\gamma$ (rad/s) & $\Omega_r$ prevista (rad/s) \\
\midrule
Decaimento da amplitude (Equação~\ref{eq:ln}) & 0,37 & 4,462 \\
Frequências $\omega_0$ e $\omega_1$ (Equação~\ref{eq:gw}) & 0,96 & 4,282 \\
\midrule
Valor experimental & --- & $4{,}21\pm0{,}04$ \\
\bottomrule
\end{tabular}
\fonte{Elaborada pelos autores.}
\end{table}

Ambas as previsões indicam diminuição de $\Omega_r$ em relação a $\omega_0$, mas ficam acima do valor experimental. O deslocamento observado ($6{,}3\%$ abaixo de $\omega_0$) é bem maior que o previsto com $\gamma=0{,}37\un{rad/s}$ ($0{,}7\%$) e mais próximo da previsão com $\gamma=0{,}96\un{rad/s}$. Isso reforça que o modelo de atrito viscoso linear, com $\omega_0$ igual ao do ar, descreve o sistema na água apenas de modo aproximado, sendo o aumento da massa efetiva do corpo pelo fluido arrastado a causa mais provável da diferença, pois reduz a frequência natural na água. Além disso, a ressonância foi identificada por inspeção visual do máximo de amplitude, o que limita a precisão de $\Omega_R$.
